\section{Weighted Least Squares and Constant Gain Learning}\label{sec:w_ls}
% =====================================================

We saw that memory, in the sense of the weight assigned to past observations, plays a key role in determining the learning outcome.
Under RLS learning with decreasing gain, agents eventually learn that returns are unpredictable.
However, in practice agents may assign more weight to recent observations, either because they believe that the data-generating process is time-varying, or because they have limited memory capacity.
To capture this, the second learning mechanism we consider is constant gain learning.
This obtains by replacing the decreasing gain $1/t$ in RLS with a constant gain $g \in (0,1)$, and as we show in Appendix \ref{app:WLS} constant gain corresponds to an on-line estimation of WLS (Weighted Least Squares) with geometrically declining weights.

\begin{equation}
    \beta_t = \beta_{t-1} + \gamma M^{-1}_{t-1} x_{t-2} (r_{t-1} - x_{t-2} \beta_{t-1}).
\end{equation}

\begin{equation}
    M_t = M_{t-1} + \gamma \left(x_{t-1}^2 - M_{t-1} \right).
\end{equation}

\begin{equation}\label{eq:constant_gain_returns}
    r_{t-1} = \eta_{t-1}
    + \log\!\Bigl(1 - \kappa e^{w \tanh(x_{t-2} \beta_{t-2} / w)}\Bigr)
    - \log\!\Bigl(1 - \kappa e^{w \tanh(x_{t-1} \beta_{t-1} / w)}\Bigr).
\end{equation} 



Given that the gain is not decreasing and that the updating depends on the stochastic term $\eta_{t-1}$, differently to RLS, the learning mechanism cannot converge to a fixed point.
However, under suitable conditions and for small values of $\gamma$, and under the same stability condition $\kappa \in (0, 1)$, it can be shown that the learning dynamics will fluctuate around the CMCE.
We formalize this in the following proposition.
\begin{proposition}[Distribution under Constant Gain Learning] \label{prop:WLS_distribution}
    Under Constant gain learning with parameter $\gamma$, and for $\kappa < 1$, for small values of $\gamma$ the estimated parameter $\beta$ approximately follows a normal distribution
\begin{equation}
    \beta_t \sim \mathcal{N} \left(\beta^*, \gamma \frac{\sigma_d^2 (1 - \kappa)}{2 \sigma^2_x}\right),
\end{equation}
where $\beta^* = 0$ is the cross-moment consistent equilibrium.
\end{proposition}
Proof. In Appendix \ref{app:proofs_WLS}.

Proposition \ref{prop:WLS_distribution} characterizes the stochastic fluctuations in beliefs under finite memory. 
The variance of the estimated coefficient, $\gamma \sigma_d^2 (2\sigma^2_x(1-\kappa))^{-1}$, is driven by three fundamental parameters.

First, the gain parameter $\gamma$ controls memory length and directly amplifies variance. Higher $\gamma$ means agents discount past data more aggressively, giving recent observations disproportionate influence. This amplifies the impact of sampling noise: a few unusual return realizations can substantially shift beliefs when historical evidence receives little weight. 

Second, dividend volatility $\sigma_d^2$ determines the fundamental uncertainty in returns. Greater return volatility generates larger forecast errors, which translate into more volatile coefficient estimates as agents attempt to fit a predictive relationship to noisy data.

Third, predictor variance $\sigma^2_x$ acts as a stabilizing force. When the predictor varies widely, each observation provides more information about the forecasting relationship, reducing estimation uncertainty. 
% =====================================================
